\documentclass[11pt]{article}
\usepackage[a4paper,margin=0.86in]{geometry}
\usepackage{amsmath,amssymb,amsthm,mathtools}
\usepackage{booktabs,array}
\usepackage{graphicx}
\usepackage{xcolor}
\usepackage{xurl}
\usepackage[hidelinks]{hyperref}
\hypersetup{
  pdftitle={Exact Memory of Direct-Sum Spectral Fan-Out},
  pdfauthor={Lluis Eriksson},
  pdfsubject={Direct-sum memory laws for passive rational-inner spectral routing}
}
\usepackage[nameinlink,capitalise]{cleveref}
\usepackage{microtype}
\usepackage{etoolbox}
\AtBeginEnvironment{thebibliography}{\scriptsize\setlength{\itemsep}{-1pt}\setlength{\parsep}{0pt}}

\newtheorem{theorem}{Theorem}[section]
\newtheorem{proposition}[theorem]{Proposition}
\newtheorem{lemma}[theorem]{Lemma}
\newtheorem{corollary}[theorem]{Corollary}
\theoremstyle{definition}
\newtheorem{definition}[theorem]{Definition}
\newtheorem{example}[theorem]{Example}
\theoremstyle{remark}
\newtheorem{remark}[theorem]{Remark}
\newtheorem{limitation}[theorem]{Limitation}

\newcommand{\C}{\mathbb C}
\newcommand{\D}{\mathbb D}
\newcommand{\T}{\mathbb T}
\newcommand{\rank}{\operatorname{rank}}
\newcommand{\ran}{\operatorname{ran}}
\newcommand{\diag}{\operatorname{diag}}
\newcommand{\Span}{\operatorname{span}}
\newcommand{\mcdeg}{\deg_{\mathrm{McM}}}
\newcommand{\nminus}{\nu_{-}}
\newcommand{\norm}[1]{\lVert#1\rVert}
\newcommand{\one}{\mathbf 1}

\title{Exact Memory of Direct-Sum Spectral Fan-Out:\\
Generic Maximality, Colliding Targets, and the Three-Line Phase Diagram}
\author{Lluis Eriksson}
\date{11 August 2026}

\begin{document}
\maketitle

\begin{abstract}
At $L$ boundary frequencies, a passive lossless network is required to send
one fixed $k$-dimensional input channel to prescribed output $k$-planes.
Previous exact fan-out laws used mutually orthogonal targets.  We show that
orthogonality is inessential and that maximal memory is generic.  If the
$L$ target planes are in direct-sum position, then every square finite
rational-inner interpolant regular at the nodes has minimum McMillan degree
\[
                         d_{\min}=k(L-1).
\]
More generally, with $r$ the dimension of their linear span,
\[
                         r-k\le d_{\min}\le k(L-1).
\]
The upper bound is constructive: all free diagonal blocks of the boundary
Pick matrix are completed by one positive leading block and a zero Schur
complement, producing a conservative realization with at most $k(L-1)$
states.  The lower bound follows from a Pick--Stein displacement and is
independent of frame gauges and node spacing.

The theorem exposes an exact-data discontinuity.  We construct $L$ direct-sum
$k$-planes whose pairwise angles tend to zero.  Their minimum degree remains
$k(L-1)$ for every nonzero opening but falls to zero at collision; the
certifying Gram margin is exactly $\epsilon^2/(1+\epsilon^2)$.  A fail-closed
noisy eigenvalue rule quantifies the precision needed to certify this cost.
For three scalar target lines we also determine the complete phase diagram:
degree zero for one line, degree one precisely for three distinct lines in a
common two-plane, and degree two otherwise.  Thus a repeated target can cost
more than three distinct targets, a rigidity invisible to span alone.
Deterministic code replays 512 arbitrary-node completions, 1,024 noisy
certificates, 20 collision fixtures, and five minimal unitary colligations,
with an independent verifier and frozen machine-readable certificate.
\end{abstract}

\noindent\textbf{Keywords:} rational inner matrix; McMillan degree;
boundary Nevanlinna--Pick interpolation; positive semidefinite completion;
Stein displacement; spectral routing; passive network; Grassmannian.

\section{The question and the answer}

A finite passive lossless device is represented here by a square rational
matrix $S$ that is analytic on the unit disk and unitary almost everywhere
on its boundary.  Its McMillan degree counts the states of a minimal
conservative realization.  Fix an isometric input frame
$X\in\C^{N\times k}$ and distinct nodes
$\zeta_1,\ldots,\zeta_L\in\T$.  At node $i$ only an output plane
$\mathcal Y_i\in\operatorname{Gr}(k,N)$ is prescribed:
\begin{equation}\label{eq:routing-task}
                 S(\zeta_i)\ran X=\mathcal Y_i,
                 \qquad i=1,\ldots,L.
\end{equation}
The interpolant is required to be regular at the nodes.  Define
\begin{equation}\label{eq:dmin}
d_{\min}(\zeta,\mathcal Y;X)
=\min\{\mcdeg S:S\text{ is square rational inner and satisfies
\eqref{eq:routing-task}}\}.
\end{equation}
The value is independent of the particular input frame: constant unitaries
move one $k$-plane to another without changing degree.

The orthogonal example treated in the preceding fan-out paper showed
$d_{\min}=k(L-1)$ for mutually orthogonal targets.  It left open whether
orthogonality creates the cost or merely makes it visible.  The answer is
the latter.

\begin{theorem}[Universal sandwich]\label{thm:sandwich}
Let $\zeta_1,\ldots,\zeta_L$ be distinct points of $\T$ and let
$\mathcal Y_1,\ldots,\mathcal Y_L$ be $k$-planes in $\C^N$.  Put
\[
                     r=\dim\Span(\mathcal Y_1,\ldots,\mathcal Y_L).
\]
Then
\begin{equation}\label{eq:sandwich}
                       r-k\le d_{\min}\le k(L-1).
\end{equation}
Both inequalities are independent of node separation.  The upper bound is
attained by an explicit positive Pick completion.
\end{theorem}

\begin{corollary}[Generic maximality]\label{cor:generic}
If the targets are in direct-sum position,
\[
                \dim\Span(\mathcal Y_1,\ldots,\mathcal Y_L)=Lk,
\]
then
\begin{equation}\label{eq:exact-law}
                            d_{\min}=k(L-1).
\end{equation}
When $N\ge Lk$, direct-sum position is open and dense in
$\operatorname{Gr}(k,N)^L$.  Hence the largest possible minimum degree is
the generic value, not an exceptional orthogonal value.
\end{corollary}

Every pair drawn from a direct-sum table has minimum degree $k$, by applying
\cref{thm:sandwich} with $L=2$.  The ratio between the global minimum and
the largest pairwise minimum is therefore $L-1$.  Pairwise feasibility does
not aggregate into global economy.

Three features distinguish the present result from a restatement of the
orthogonal construction.  First, the target Gram matrix may be fully dense.
Second, the upper bound now requires a block completion rather than a scalar
Cauchy shift.  Third, direct-sum targets may be arbitrarily close in the
Grassmannian, producing the discontinuity proved in \cref{sec:collision}.

\section{The Pick--Stein interface}\label{sec:interface}

Choose arbitrary isometric frames $W_i\in\C^{N\times k}$ with
$\ran W_i=\mathcal Y_i$, and form the calibrated block Gram matrix
\begin{equation}\label{eq:gram}
 G=[W_i^*W_j]_{i,j=1}^L=W^*W,
 \qquad W=[W_1\ \cdots\ W_L].
\end{equation}
Changing frames inside the planes conjugates $G$ by a block-diagonal unitary.
In particular,
\begin{equation}\label{eq:gram-rank}
                      \rank G=\rank W=r
\end{equation}
is intrinsic.

Let
\[
 Z=\diag(\zeta_1,\ldots,\zeta_L)\otimes I_k,
 \qquad J_L=\one\one^*.
\]
The boundary Pick variables are the blocks of a Hermitian matrix
$P=[P_{ij}]\in\C^{Lk\times Lk}$.  Its off-diagonal blocks are fixed:
\begin{equation}\label{eq:pick-offdiag}
 P_{ij}=\frac{I_k-W_i^*W_j}{1-\overline{\zeta_i}\zeta_j},
 \qquad i\ne j.
\end{equation}
The diagonal blocks are free because both numerator and denominator vanish
on the boundary.  Equivalently, every completion satisfies the displacement
identity
\begin{equation}\label{eq:stein}
                 P-Z^*PZ=J_L\otimes I_k-G.
\end{equation}

We isolate the standard interpolation bridge in the precise form used below.

\begin{proposition}[Rank--degree bridge]\label{prop:bridge}
If a matrix $P\succeq0$ satisfies \eqref{eq:stein}, then there is a square
rational-inner $S$, regular at the prescribed nodes after removal of
decoupled unit-circle modes, such that
$S(\zeta_i)X=W_i$ and
\[
                           \mcdeg S\le\rank P.
\]
Conversely, an interpolant of degree $d$ produces such a positive $P$ with
$\rank P\le d$.
\end{proposition}

\begin{proof}
Factor $P=F^*F$ with $F=[F_1\ \cdots\ F_L]$ and
$F_i\in\C^{q\times k}$, where $q=\rank P$.  Identity \eqref{eq:stein}
says that the two block-column families
\[
  \begin{bmatrix}\zeta_iF_i\\X\end{bmatrix}
  \quad\text{and}\quad
  \begin{bmatrix}F_i\\W_i\end{bmatrix}
\]
have the same Gram matrix.  The induced isometry extends to a unitary
colligation
\[
 U=\begin{bmatrix}A&B\\C&D\end{bmatrix}
 \quad\text{on }\C^q\oplus\C^N.
\]
Its transfer matrix
\[
                    S(z)=D+zC(I-zA)^{-1}B
\]
is rational inner and has degree at most $q$.  The interpolation equations
give $(I-\zeta_iA)F_i=BX$ and then $S(\zeta_i)X=W_i$.  If the chosen unitary
extension contains a unit-circle state mode, contractivity forces that mode
to be reducing and decoupled from $B$ and $C$; deleting it leaves the
transfer matrix and the interpolation data unchanged and makes all nodes
regular.

Conversely, the de Branges--Rovnyak kernel of a finite rational-inner matrix,
\[
 K_S(z,w)=\frac{I-S(z)^*S(w)}{1-\bar z w},
\]
has a finite Kolmogorov factor of dimension $d=\mcdeg S$.  Taking regular
boundary limits and compressing by $X$ gives $P\succeq0$ with
$\rank P\le d$, off-diagonal blocks \eqref{eq:pick-offdiag}, and hence
\eqref{eq:stein}.  This is the finite-dimensional lurking-isometry form of
matrix boundary Nevanlinna--Pick interpolation
\cite{BolotnikovDym2006,BallKim1994,GombaniMichaletzky2007}.
\end{proof}

The proposition separates the new geometry from established realization
machinery.  Our contribution is not the Schur complement or the
rank--degree bridge individually, but the sharp completion and lower bound
that make them meet for generic subspace data.

\section{A universal rank-\texorpdfstring{$k(L-1)$}{k(L-1)} completion}\label{sec:upper}

Partition the partially specified matrix in \eqref{eq:pick-offdiag} after its
first $L-1$ block rows and columns.  Let $A_0$ be the Hermitian
$k(L-1)\times k(L-1)$ matrix whose off-diagonal blocks are the corresponding
blocks in \eqref{eq:pick-offdiag} and whose diagonal blocks are zero.  Let
$B_0$ stack the prescribed blocks $P_{1L},\ldots,P_{L-1,L}$.

\begin{lemma}[Zero-Schur completion]\label{lem:completion}
For every $t>\norm{A_0}$, set
\begin{equation}\label{eq:completion}
 Q=tI+A_0,
 \qquad
 P=\begin{bmatrix}
 Q&B_0\\ B_0^*&B_0^*Q^{-1}B_0
 \end{bmatrix}.
\end{equation}
Then $P\succeq0$, it has all the prescribed off-diagonal blocks, and
\[
                         \rank P=k(L-1).
\]
\end{lemma}

\begin{proof}
The choice of $t$ gives $Q\succ0$.  The Schur complement of $Q$ in $P$ is
zero, so $P\succeq0$ and $\rank P=\rank Q=k(L-1)$.  Only diagonal blocks
were selected freely: the leading diagonal blocks are those of $tI$, and
the final diagonal block is $B_0^*Q^{-1}B_0$.  Thus every off-diagonal block
remains exactly \eqref{eq:pick-offdiag}.
\end{proof}

Positive semidefinite completion by a zero Schur complement is elementary
and belongs to the general matrix-completion toolkit
\cite{Dancis1992,LaurentVarvitsiotis2014}.  Its consequence here is less
obvious: arbitrary boundary subspace data always admit a passive lossless
interpolant with a state count linear in both $k$ and $L$.

\begin{proof}[Proof of the upper bound in \cref{thm:sandwich}]
Apply \cref{lem:completion}.  Its matrix obeys \eqref{eq:stein}; the diagonal
blocks require no check because $|\zeta_i|=1$ and $W_i^*W_i=I_k$.
\Cref{prop:bridge} produces an interpolant of degree at most $k(L-1)$.
\end{proof}

The construction is algorithmic.  It needs one Hermitian eigenvalue bound,
one positive linear solve, a factorization of $P$, and a unitary completion
of the lurking isometry.  It does not optimize conditioning: increasing $t$
preserves rank but can enlarge stored-energy scales.  Degree and numerical
conditioning are different resources.

\section{The span obstruction and exact generic law}\label{sec:lower}

We use two elementary inertia facts.  If $P\succeq0$ has rank $q$, then
\begin{equation}\label{eq:displacement-inertia}
                      \nminus(P-Z^*PZ)\le q,
\end{equation}
because the negative inertia of a difference $P-R$ cannot exceed
$\rank R$, here $R=Z^*PZ$.  Also, adding a positive matrix of rank at most
$k$ can remove at most $k$ negative directions.  Since $G\succeq0$,
\begin{equation}\label{eq:rank-inertia}
 \nminus(J_L\otimes I_k-G)
 \ge \nminus(-G)-k
 =\rank G-k.
\end{equation}

\begin{proof}[Proof of the lower bound in \cref{thm:sandwich}]
Let $S(z)=D+zC(I-zA)^{-1}B$ be a minimal realization of degree $d$.
The resolvent identity gives
\[
 S(z)-S(w)=(z-w)C(I-zA)^{-1}(I-wA)^{-1}B.
\]
Hence every column of $(S(\zeta_i)-S(\zeta_1))X$ lies in the common space
$\ran C$, of dimension at most $d$.  Therefore
\[
 \mathcal Y_1+\cdots+\mathcal Y_L
 \subseteq \mathcal Y_1+\ran C,
 \qquad r\le k+d.
\]
This proves $d\ge r-k$ directly.  Equivalently, the Pick matrix from
\cref{prop:bridge} obeys
$\nminus(P-Z^*PZ)\le\rank P\le d$, while
\eqref{eq:rank-inertia} gives
$\nminus(P-Z^*PZ)\ge r-k$.  The two arguments audit the same obstruction
from realization and interpolation viewpoints.
\end{proof}

\begin{proof}[Proof of \cref{cor:generic}]
Direct-sum position gives $r=Lk$, so the two sides of
\eqref{eq:sandwich} coincide.  Full column rank of the concatenated frame
$W$ is an open condition.  Its failure is the zero set of all $Lk$-minors;
when $N\ge Lk$, at least one such minor is nonzero for an orthogonal table,
so the full-rank locus is dense.
\end{proof}

The conclusion is stronger than ``most targets need many states.''  It
identifies the exact maximum and proves that this maximum occupies an open
dense set.  Orthogonal targets supply a transparent point in that set, not
the mechanism behind the law.

\section{Colliding targets and exact-data discontinuity}\label{sec:collision}

Let $E_0,E_1,\ldots,E_L$ be mutually orthogonal isometric $k$-frames in
$\C^{(L+1)k}$.  For $\epsilon\ge0$ define
\begin{equation}\label{eq:collision-family}
 W_i(\epsilon)=\frac{E_0+\epsilon E_i}{\sqrt{1+\epsilon^2}},
 \qquad \mathcal Y_i(\epsilon)=\ran W_i(\epsilon).
\end{equation}
For distinct $i,j$,
\begin{equation}\label{eq:collision-overlap}
 W_i(\epsilon)^*W_j(\epsilon)=\frac{1}{1+\epsilon^2}I_k.
\end{equation}
Thus all principal angles tend to zero as $\epsilon\downarrow0$.

\begin{corollary}[Maximal memory at arbitrarily small angles]
\label{cor:collision}
For the family \eqref{eq:collision-family},
\begin{equation}\label{eq:collision-degree}
 d_{\min}(\epsilon)=
 \begin{cases}
 k(L-1),&\epsilon>0,\\
 0,&\epsilon=0.
 \end{cases}
\end{equation}
The block Gram spectrum is
\begin{equation}\label{eq:collision-spectrum}
 \left\{\frac{L+\epsilon^2}{1+\epsilon^2}\right\}^{k}
 \cup
 \left\{\frac{\epsilon^2}{1+\epsilon^2}\right\}^{k(L-1)}.
\end{equation}
\end{corollary}

\begin{proof}
The private $E_i$ components make the $Lk$ columns of
$[W_1(\epsilon)\ \cdots\ W_L(\epsilon)]$ independent for every
$\epsilon>0$.  Apply \cref{cor:generic}.  At zero all targets equal
$\ran E_0$, and a constant unitary realizes the table.  Finally,
\[
 G_\epsilon=\frac{J_L+\epsilon^2I_L}{1+\epsilon^2}\otimes I_k,
\]
whose spectrum is \eqref{eq:collision-spectrum}.
\end{proof}

\begin{figure}[t]
\centering
\includegraphics[width=0.94\linewidth]{figures/generic_fanout_memory.pdf}
\caption{Left: six target lines approach coincidence while the exact minimum
degree remains five for every positive opening.  The Gram margin closes
quadratically.  Right: every pair costs one state, whereas the generic global
task costs $L-1$.  Curves show exact formulas, not fitted asymptotics.}
\label{fig:collision}
\end{figure}

There is no contradiction with continuity of approximate interpolation.
The feasible set changes discontinuously only because equality at every node
is exact.  If a tolerance is admitted, a constant device approximates
\eqref{eq:collision-family} increasingly well.  The small Gram eigenvalue in
\eqref{eq:collision-spectrum} is precisely the statistical price of
distinguishing an expensive positive opening from a free collision.

\section{A gauge-invariant finite-error certificate}\label{sec:noisy}

Let $\Pi_i=W_iW_i^*$ be the orthogonal projector onto $\mathcal Y_i$ and set
\begin{equation}\label{eq:projector-sum}
                           H=\sum_{i=1}^L\Pi_i.
\end{equation}
This observable is independent of all frame gauges, and
$\rank H=\dim\sum_i\mathcal Y_i$.  Suppose projector estimates satisfy
$\norm{\widehat\Pi_i-\Pi_i}\le\epsilon_i$.  With
$\widehat H=\sum_i\widehat\Pi_i$ and $\eta=\sum_i\epsilon_i$, one has
\begin{equation}\label{eq:error-model}
                         \norm{\widehat H-H}\le\eta.
\end{equation}
The error bound may come from deterministic calibration guarantees or a
high-probability confidence event.  No distributional assumption is needed
inside the theorem.

\begin{theorem}[Noisy span certificate]\label{thm:noisy}
Let
\[
 q_\eta=\#\{j:\lambda_j(\widehat H)>\eta\}.
\]
Every exact rational-inner interpolant satisfies
\begin{equation}\label{eq:noisy-bound}
                        \mcdeg S\ge\max\{0,q_\eta-k\}.
\end{equation}
In particular, if $q_\eta=Lk$, then $d_{\min}=k(L-1)$.
\end{theorem}

\begin{proof}
Weyl's inequality and \eqref{eq:error-model} imply that each eigenvalue of
$\widehat H$ above $\eta$ corresponds to a strictly positive eigenvalue of
$H$.  Hence $\rank H\ge q_\eta$.  Apply the lower half of
\cref{thm:sandwich}; when $q_\eta=Lk$, also apply its universal upper half.
\end{proof}

The rule is fail-closed: a large $\eta$ can make it abstain but cannot make it
overstate the span on the event \eqref{eq:error-model}.  For
\eqref{eq:collision-family}, certification requires an error smaller than a
quantity of order $\epsilon^2$.  Nearly coincident planes are geometrically
close but quadratically hard to certify as algebraically independent.

Unlike a block Gram estimate, \eqref{eq:projector-sum} needs no relative
$U(k)$ calibration between frames.  A significantly negative eigenvalue
$\lambda_{\min}(\widehat H)<-\eta$ falsifies the declared projector/error
model, because the exact $H$ is positive semidefinite.

\section{The exact phase diagram for three target lines}\label{sec:three}

The span lower bound is sharp on the open direct-sum stratum, but not on every
singular stratum.  The smallest nontrivial example already contains a
rigidity phenomenon.  Let $L=3$ and $k=1$.

\begin{theorem}[Three-line phase diagram]\label{thm:three-lines}
For three target lines $\mathcal Y_1,\mathcal Y_2,\mathcal Y_3$ at distinct
boundary nodes,
\begin{equation}\label{eq:three-lines}
d_{\min}=
\begin{cases}
0,&\mathcal Y_1=\mathcal Y_2=\mathcal Y_3,\\
1,&\text{the three lines are pairwise distinct and span a two-plane},\\
2,&\text{otherwise}.
\end{cases}
\end{equation}
\end{theorem}

\begin{proof}
A constant unitary realizes, and is forced by, the first case.  A degree-one
rational-inner matrix is, up to constant unitary factors, one rank-one
Blaschke--Potapov factor
\[
                     \diag(1,b(z),I),
\]
where $b$ is a disk automorphism.  Applied to a fixed input line, its output
is either constant or has the form
\begin{equation}\label{eq:potapov-orbit}
                        [u+\rho b(z)v],
\end{equation}
for orthogonal vectors $u,v$ and $\rho>0$.  The nonconstant orbit is
injective on $\T$ and lies in $\Span\{u,v\}$.  Thus degree one cannot produce
a repeated line followed by a different line, nor three lines spanning
dimension three.

Conversely, identify the common two-plane of three distinct targets with
$\C^2$.  Under the Bloch-sphere representation the three points determine a
circle.  A constant $U(2)$ rotation sends that circle to a latitude, so the
three lines acquire the form $[u+\rho e^{i\theta_i}v]$.  Choose the circle
orientation so that $(e^{i\theta_1},e^{i\theta_2},e^{i\theta_3})$ and
$(\zeta_1,\zeta_2,\zeta_3)$ have the same cyclic order.  The disk automorphism
group acts triply transitively on oriented boundary triples, hence some
degree-one Blaschke factor satisfies $b(\zeta_i)=e^{i\theta_i}$.  Extending
the two constant unitaries to $\C^N$ realizes \eqref{eq:potapov-orbit}.

If the lines span dimension three, \cref{thm:sandwich} gives degree at least
two.  If exactly two distinct lines occur, injectivity of every nonconstant
degree-one orbit gives the same lower bound.  The universal upper bound with
$L=3,k=1$ supplies degree at most two in both cases.
\end{proof}

The second obstruction is genuinely global but not a rank obstruction.
Three distinct coplanar lines cost one state; replacing one by a repeated
line can raise the cost to two.  Algebraic incidence and interpolation order
therefore matter on the rank-deficient boundary of the generic stratum.
For four or more scalar nodes, degree-one feasibility additionally imposes
Möbius cross-ratio compatibility; the three-node case is special because an
oriented boundary triple has no remaining modulus.

\section{Reproducible synthesis and stress tests}\label{sec:computation}

The proof is exact and does not depend on floating-point evidence.  The
computation checks that the proposed completion and compiler behave as
claimed away from the orthogonal case.  The producer is
\href{https://github.com/lluiseriksson/finite-sample-spectral-certificates/blob/research/generic-fanout-memory/research/generic_fanout_memory_certificate.py}{\texttt{research/generic\_fanout\_memory\_certificate.py}};
the independent implementation is
\href{https://github.com/lluiseriksson/finite-sample-spectral-certificates/blob/research/generic-fanout-memory/verification/verify_generic_fanout_memory.py}{\texttt{verification/verify\_generic\_fanout\_memory.py}}.
The verifier imports no producer functions.

The deterministic campaign contains four blocks:
\begin{enumerate}
\item 512 random arbitrary-node, nonorthogonal target tables with
      $3\le L\le8$ and $1\le k\le3$;
\item 20 near-collision fixtures down to $\epsilon=10^{-3}$;
\item five unitary colligations, checked for interpolation, unitarity,
      controllability, observability, and the predicted state dimension;
\item 1,024 noisy projector-sum trials, half direct-sum and half rank deficient;
\item four clustered-node fixtures with gaps from $10^{-2}$ to $10^{-8}$.
\end{enumerate}

\begin{table}[t]
\centering
\caption{Representative colligations synthesized from dense target Gram
matrices.  Errors are spectral norms.}
\label{tab:colligations}
\small
\begin{tabular}{@{}rrrrrr@{}}
\toprule
$L$ & $k$ & states & interpolation & colligation unitary & $\rho(A)$\\
\midrule
3&1&2&$1.72\,10^{-14}$&$1.26\,10^{-14}$&0.7881\\
4&1&3&$2.16\,10^{-14}$&$3.15\,10^{-14}$&0.9610\\
3&2&4&$3.05\,10^{-13}$&$1.87\,10^{-13}$&0.9918\\
4&2&6&$3.28\,10^{-14}$&$3.01\,10^{-14}$&0.8677\\
5&2&8&$2.14\,10^{-14}$&$6.44\,10^{-14}$&0.8584\\
\bottomrule
\end{tabular}
\end{table}

Across the 512 completion trials, the worst relative Stein residual was
$4.23\times10^{-16}$ and the worst relative violation of positive
semidefiniteness was $5.06\times10^{-16}$.  Every completion had rank
$k(L-1)$.  All 512 direct-sum noisy fixtures were certified and none of the
1,024 trials overcertified the true span bound.  These counts and exact
floating-point outputs are frozen in a JSON certificate rather than copied
from plotting code.

Close nodes are an exact nonissue but a numerical stressor: the entries of
the Pick completion scale inversely with the node gap.  The clustered-node
fixtures retain relative Stein residuals below the declared tolerance down to
$10^{-8}$, while recording the floating-point rank separately from the exact
rank guaranteed by the zero Schur complement.  We do not claim reliable
double-precision colligation synthesis at the smallest gaps.

The release gate performs three distinct checks: hash verification of frozen
artifacts, numerical replay in a temporary directory, and an independent
formula-level audit.  The immutable bundle is linked at
\href{https://github.com/lluiseriksson/finite-sample-spectral-certificates/releases/tag/v2.2-generic-fanout-memory}{\texttt{v2.2-generic-fanout-memory}}.

\section{Scope, relation to prior work, and limitations}\label{sec:scope}

Boundary matrix Schur interpolation, minimal rational interpolation, and
degree control have a substantial literature
\cite{BolotnikovDym2006,AntoulasEtAl1990,BallKim1994,
GombaniMichaletzky2007,BallBolotnikov2018}.  Subspace interpolation has also
been used to design matrix all-pass filters
\cite{RapisardaWillems1997,BharathEtAl2023}.  Positive semidefinite completion
and rank questions are classical in their own right
\cite{Dancis1992,LaurentVarvitsiotis2014}.  We do not claim any of these
general mechanisms as new.

The contribution is the following conjunction for passive spectral routing:
the universal sharp sandwich \eqref{eq:sandwich}; the exact open-dense law
\eqref{eq:exact-law}; a single explicit completion that proves the upper
bound for arbitrary subspaces and nodes; the collision discontinuity
\eqref{eq:collision-degree}; its finite-error certification threshold; and
the complete singular-stratum classification \eqref{eq:three-lines}.
Compared with the orthogonal fan-out theorem, the result replaces symmetry
by a generic geometric condition and identifies what fails when that
condition degenerates.

Several boundaries should be kept explicit.
\begin{itemize}
\item The task is exact interpolation at finitely many distinct boundary
      nodes.  Approximate, broadband, or derivative constraints define
      different optimization problems.
\item The transfer is finite-dimensional, square, rational, and inner.
      Loss, gain, infinite delays, and nonrational distributed systems are
      outside the theorem.
\item McMillan degree counts dynamical states.  It does not by itself bound
      dwell time, bandwidth, component sensitivity, or implementation loss.
\item The noisy theorem is conditional on valid operator-norm error bounds
      for the estimated projectors.  The simulations illustrate
      the logic; they are not a substitute for an experimental noise model.
\item Outside direct-sum position, $r-k$ need not be exact.  The repeated-line case in
      \cref{thm:three-lines} is the smallest counterexample.  A complete
      stratification for general $L$ and $k$ remains open.
\end{itemize}

\section{Conclusion}

Generic spectral fan-out is maximally expensive.  The exact state count
$k(L-1)$ does not require orthogonal destinations, large angles, or well
separated frequencies.  It follows from the algebraic independence of the
target planes, and it persists at arbitrarily small positive angles.  The
universal completion proves that no routing table costs more; the span
inertia proves that a direct-sum table costs no less.

The collision family separates exact complexity from geometric distance.
Its quadratic Gram margin also turns that separation into an operational
statement about certification precision.  Finally, the three-line phase
diagram shows that span is the complete story only on the generic stratum:
at singular configurations, Möbius rigidity supplies an additional memory
obstruction.  Together these results replace one symmetric example by a
global law, a sharp boundary phenomenon, and a reproducible compiler.

\paragraph{Data and code availability.}
Source, producer, independent verifier, frozen certificate, figure, build
logs, and artifact manifest are contained in the repository and immutable
release linked in \cref{sec:computation}.  The numerical evidence is
supporting material; all mathematical claims are proved in the text.

\bibliographystyle{plain}
\bibliography{references}

\end{document}
